% Archivo: gomez_hernandez_marco_tfg.tex
\documentclass[epsbased,copyright,final,printable,covers,tfg]{tfgtfmthesisuam}

% --- Metadatos del Estudiante y TFG ---
\advisor{Nombre del Tutor}
\author{Marco Gomez Hernandez}
\title{Título del Trabajo Fin de Grado}
\levelin{Ingeniería Informática}
\copyrightdate{Junio 2026}

% --- Datos de la Facultad (EPS) ---
\coverdata{
  Escuela Politécnica Superior \\
  Universidad Autónoma de Madrid \\
  C\textbackslash Francisco Tomás y Valiente nº 11
}

% --- Archivos de Inicio (Resúmenes y Agradecimientos) ---
\prefacefile{inicio/prefacio}       % Opcional
\ackfile{inicio/agradecimientos}    % Obligatorio
\resumenfile{inicio/resumen}        % Obligatorio
\abstractfile{inicio/abstract}      % Obligatorio

\keywords{Palabra clave 1, Palabra clave 2, Palabra clave 3}
\palabrasclave{Keyword 1, Keyword 2, Keyword 3}

% --- Configuración de fuentes de datos ---
\bibliographyconfig{bib/bibliografia}
\graphicsdir{img}
\logosdir{img}
\codesdir{codes}    % Carpeta para códigos fuente
\datadir{data}      % Carpeta para datos de gráficas

\begin{document}

% NOTA: Los índices (Figuras, Tablas, Ecuaciones, Algoritmos, Códigos) 
% se generan AUTOMÁTICAMENTE al compilar gracias a la clase tfgtfmthesisuam.

% --- CAPÍTULOS ---
% La sintaxis de esta plantilla es: \chapter{Título del Capítulo}{ruta/del/fichero}

% 1. Introducción (Motivación, Objetivos, Estructura)
\chapter{Introducción\label{CAP:INTRO}}{capitulos/01_Introduccion}

% 2. Estado del Arte / Conceptos Teóricos
\chapter{Estado del Arte\label{CAP:ESTADO}}{capitulos/02_EstadoArte}

% 3. Desarrollo / Metodología (Aquí ponemos ejemplos para rellenar los índices)
\chapter{Desarrollo y Ejemplos\label{CAP:DESARROLLO}}{capitulos/03_Desarrollo}

% 4. Conclusiones y Trabajo Futuro
\chapter{Conclusiones\label{CAP:CONCLUSIONES}}{capitulos/04_Conclusiones}

% --- APÉNDICES ---
\appendix
\chapter{Manual de Usuario\label{APD:MANUAL}}{apendices/apendiceA}

\end{document}